\section{Resolvent and determinant coefficients}
\begin{proposition}[The resolvent]\label{T10:prop:resolvent}
For real $z$ with $|z|>1$,
\begin{equation}\label{T10:eq:resolvent}
 \varkappa\Bigl(\frac\lambda{\lambda-z}\Bigr)
 =\frac{4z^2}{\pi^2(z^2-1)\sqrt{2z^2-1}}
   \operatorname{artanh}\frac1{\sqrt{2z^2-1}} .
\end{equation}
\end{proposition}

\begin{proof}
Since $\varkappa$ annihilates constants,
$\varkappa(\lambda/(\lambda-z))=z\varkappa(1/(\lambda-z))$.
Set $B=2z^2-1>1$ and $v=2\lambda^2-1$ in
\eqref{T10:eq:density}. Then
$\pi^2\varpi(\lambda)\,d\lambda=dv/[(1-v)\sqrt v]$ and
\[
 \frac1{\lambda-z}+\frac1{-\lambda-z}-\frac1{1-z}-\frac1{-1-z}
 =\frac{4z(1-v)}{(B-v)(B-1)}.
\]
The endpoint factor cancels, so
\[
 \pi^2\varkappa\Bigl(\frac1{\lambda-z}\Bigr)
 =\frac{4z}{B-1}\int_0^1\frac{dv}{(B-v)\sqrt v}
 =\frac{8z}{(B-1)\sqrt B}\operatorname{artanh}(B^{-1/2}).
\]
Multiply by $z$ and use $B-1=2(z^2-1)$.
\end{proof}

Since $\mathcal K_\tau(\mathcal K_\tau-z)^{-1}$ is trace class for $|z|>1$ by
Proposition~\ref{prop:traces}, \eqref{T10:eq:resolvent} specializes
Theorem~\ref{T10:conj:law} to a resolvent trace formula: for each fixed real
$z$ with $|z|>1$ the theorem gives
$\operatorname{tr}[\mathcal K_\tau(\mathcal K_\tau-z)^{-1}]
=4\tau(1-z^2)^{-1}+\varkappa(\lambda/(\lambda-z))\log\tau+O(1)$.  Expanding
\eqref{T10:eq:resolvent} at $z=\infty$ returns
$2\pi^{-2}z^{-2}+O(z^{-4})$, consistent with
$\pi^2\varkappa(\lambda^2)=-2$; the two branch points at $|z|=2^{-1/2}$ and
$|z|=1$ are the two edges of the support of $\varpi$.

\begin{corollary}[A Fredholm determinant law]\label{v7:cor:determinant}
For real $|z|<1$ and the phase-shifted operator $\mathcal K_{\tau,\phi}$,
with kernel $[\sin(2\pi\tau x+\phi)-\sin(2\pi\tau y+\phi)]/[\pi(x-y)]$ and
treated in Proposition~\ref{v6:prop:allphase}, as $\tau\to\infty$,
\begin{equation}\label{v7:eq:determinant}
 \log\det(I-z\mathcal K_{\tau,\phi})
 =2\tau\log(1-z^2)
 +\frac2{\pi^2}\operatorname{artanh}^2
       \!\left(\frac{z}{\sqrt{2-z^2}}\right)\log\tau+O_z(1).
\end{equation}
The remainder is uniform in $\phi$ and locally uniform in $z\in(-1,1)$.
\end{corollary}
\begin{proof}
For $f_z(x)=\log(1-zx)$, all factors $1-z\lambda_j$ are positive and
$\sum_j|f_z(\lambda_j)|<\infty$.  Thus the real Fredholm logarithm is
$\tr f_z(\mathcal K_{\tau,\phi})$.  The test vanishes at zero and its
first two derivatives are uniformly bounded when $z$ ranges over a
compact subinterval of $(-1,1)$. Let $\beta(z)=\varkappa(f_z)$.
For $z\ne0$, differentiation of the convergent coefficient integral
and Proposition~\ref{T10:prop:resolvent} give
\[
 \beta'(z)=z^{-1}\varkappa\!\left(\frac{x}{x-1/z}\right)
 =\frac{4\operatorname{artanh}(z/\sqrt{2-z^2})}
 {\pi^2(1-z^2)\sqrt{2-z^2}}.
\]
This is the derivative of
$2\pi^{-2}\operatorname{artanh}^2(z/\sqrt{2-z^2})$.
Both functions vanish at zero, so integration identifies the coefficient;
the uniform phase trace law proves \eqref{v7:eq:determinant}.
\end{proof}

\begin{remark}[The square and a comparison]\label{rem:det-square}
For $z\in\mathbb C\setminus(( -\infty,-1]\cup[1,\infty))$, define
$\log\det(I-z\mathcal K_{\tau,\phi})$ as the holomorphic logarithm
normalized to zero at $z=0$, equivalently the absolutely convergent sum
$\sum_j\operatorname{Log}(1-z\lambda_j)$, where $\lambda_j$ are the eigenvalues of $\mathcal K_{\tau,\phi}$ and
$\operatorname{Log}$ is the scalar principal logarithm.  Applying the uniform phase trace law
to the real and imaginary parts of this scalar test extends
Corollary~\ref{v7:cor:determinant} to that domain.
Likewise, for $w\in\mathbb C\setminus(-\infty,-1]$, put
$\log\det(I+w\mathcal K_{\tau,\phi}^2)
=\sum_j\operatorname{Log}(1+w\lambda_j^2)$.  Then
\begin{equation}\label{eq:det-square}
 \log\det(I+w\mathcal K_{\tau,\phi}^2)
 =4\tau\operatorname{Log}(1+w)
  -\frac4{\pi^2}\arcsin^2\!\Bigl(\tfrac12\sqrt{\tfrac{2w}{1+w}}\Bigr)\log\tau
  +O_w(1).
\end{equation}
The squared inverse-sine expression denotes its analytic continuation
from $w>0$ throughout the stated slit plane, with value zero at $w=0$.
For $w>0$, add the two logarithms at $z=\pm i\sqrt w$ and use
$\operatorname{artanh}(iy)=i\arctan y$ and
$\arctan\sqrt{w/(2+w)}=\arcsin\bigl(\tfrac12\sqrt{2w/(1+w)}\bigr)$.
Analytic continuation of the coefficient integral then gives the formula
on the slit plane; the trace-law error bound applies directly at each $w$
and is uniform in $\phi$.  At $w=1$ the coefficient is $-\tfrac19$, so
$\log\det(I+\mathcal K_{\tau,\phi}^2)=4\tau\log2-\tfrac19\log\tau+O(1)$.

For comparison, Gebert and K\"uttler's generalized Hilbert matrix has
entries $(H_N^\delta)_{jk}=\sin((1+\delta)\pi)/[\pi(j+k+1+\delta)]$,
$0\le j,k<N$.  For $\delta<-1/2$ outside the negative integers and
half-integers, their Corollary~1.3 gives
$\log\det(I-(H_N^\delta)^2)=-\delta^2\log N+O_\delta(1)$
\cite{GebertKuettler2026}.  This is an analogous determinant law for a
different signed family; their critical half-integers have a weaker
remainder, as specified in their Remark~1.4.
\end{remark}

These formulas give bounded remainders. Theorem~\ref{v11:thm:phase-constant} proves their convergence along fixed phases by localizing the physical boundary corrections.  For the unsquared determinant with real $|z|<1$, the three-dimensional model
symbol $a(\theta)$ of Section~\ref{app:general-trace} has a coefficient of
$\log\tau$ equal to the exponent $-\operatorname{tr}\mathfrak L^2$,
$\mathfrak L=(2\pi i)^{-1}\log[(I-zX_{1,0})^{-1}(I-zX_{0,-1})]$, of the block
Fisher--Hartwig asymptotics \cite[Corollary~2.4]{BasorEhrhardtVirtanen2024},
which give a convergent constant term for the block Toeplitz truncation of
$I-za$ itself.


